\documentclass[a4paper,12pt]{article}

%% -------
%% List of packages must be dynamic
%% -------
\usepackage[utf8]{inputenc}
\usepackage[T1]{fontenc}
\usepackage[ngerman]{babel} 
\usepackage{amsmath, amssymb} 
\usepackage{graphicx} 
\usepackage{subcaption}
\usepackage{hyperref} 
\usepackage{geometry} 
\usepackage{array}
\usepackage{enumitem}
\usepackage{titlesec}
\usepackage{listings}
\usepackage{tikz} \usetikzlibrary{decorations.pathreplacing,arrows.meta,positioning}
\usepackage{listings}
\usepackage{tabularx}
\usepackage[most]{tcolorbox}
\usepackage{comment}

%% -------
%% General configuration
%% -------
\usetikzlibrary{arrows, automata, positioning}
\geometry{a4paper, left=2cm, right=2cm, top=2cm, bottom=2cm}
\def \runninghead {Running head goes here ... }

%% -------
%% Header and Footer
%% -------
\usepackage{fancyhdr}
\renewcommand{\headrulewidth}{0.4pt}  % Dicke des Striches
\renewcommand{\footrulewidth}{0.4pt}  % Dicke des Striches
\pagestyle{fancy}
\lhead{Matrikelnr.:}
\chead{}
\rhead{}
\lfoot{\runninghead}
\cfoot{}
\rfoot{Seite \thepage}

%% -------
%% Some more configuration
%% -------
\title{\textbf{\runninghead}} 
\author{}
\date{04.02.2025} 

%% change subsections to have no linebreak afterwards
\renewcommand{\thesubsection}{\alph{subsection})}
\titleformat{\subsection}[runin]  % Run-in format
  {\normalfont\bfseries}          % Style (bold text)
  {\thesubsection}{1em}{}         % Numbering with spacing

%% remove indents at start of new paragraphs
\setlength{\parindent}{0pt}

%% -------
%% Create a box for displaying solutions
%% -------

% Define a flag to show or hide solution
\newif\ifshowsolutions

% Change to \showsolutionstrue to show solutions
% Change to \showsolutionsfalse to hide solutions
\showsolutionstrue

% Define the solution environment if necessary
\ifshowsolutions
  \newtcolorbox{solution}{
    colback=red!80,        % Darker red background
    colframe=red!90!black, % Border color
    fontupper=\color{white}\footnotesize, % Make text white and small
    title=Solution,
    boxrule=0.8pt,
    arc=4pt,
    top=6pt,
    bottom=6pt,
    left=6pt,
    right=6pt
  }
\else %% hides the solution environment
  \excludecomment{solution}
\fi

\begin{document}

%% -------
%% Front page start
%% front pages are individual per exam
%% -------
\thispagestyle{empty} %% no header/footer for first page
\begin{center}
    \textbf{\LARGE Exam title goes here} \\
    \vspace{0.5cm}
    \textbf{Sommersemester 25 -- 2. Termin -- 17.10.2025} \\
    \vspace{0.2cm}
    \textbf{Prof. Dr. Marc Hesenius}
\end{center}

\vspace{0.5cm}
\noindent
%\textbf{Vor- und Nachname:}
%\underline{\hspace{6cm}} \hfill 
\textbf{Matrikelnr.:} \underline{\hspace{8cm}} \\

\begin{itemize}
    \item Bitte tragen Sie vor Beginn der Bearbeitung auf allen Blättern Ihren \textbf{Namen} und Ihre \textbf{Matrikelnummer} ein. Sie erhalten dafür 5 Minuten extra.
    \item Die Bearbeitungszeit beträgt \textbf{120 Minuten}.
    \item Legen Sie Ihren Studentenausweis oder einen amtlichen \textbf{Lichtbildausweis} (Personalausweis oder Reisepass) bereit.
    \item Die \textbf{Heftung} der Klausurseiten muss bestehen bleiben!
    \item Falls der gegebene Platz für die Lösungen nicht ausreichen sollte, benutzen Sie bitte die Rückseiten. Verweisen Sie an geeigneter Stelle mit Angabe der Seitennummer auf die Fortsetzung der Aufgabe.
    \item Es sind ausschließlich dokumentenechte Stifte zulässig. \textbf{Keine Bleistifte! Keine Rotstifte! Kein Tipp-Ex!}
    \item Ein beidseitig beschriebenes oder bedrucktes DIN-A4-Blatt ist das einzige erlaubte Hilfsmittel. 
    \item Auf dem Tisch befinden sich ausschließlich: Stifte, Studentenausweis, Lichtbildausweis und das DIN-A4-Blatt.
    \item Weitere persönliche Gegenstände verbleiben in den Taschen. Mobiltelefone, Taschenrechner, \textbf{Smartwatches} und andere elektronische Geräte werden ausgeschaltet vom Arbeitsplatz entfernt.
\end{itemize}

%% table should be dynamic depending on number of chosen topics
		\vspace{0.5cm}
		\hrule
	
		\begin{tabularx}{\linewidth}{|p{2.05cm}||p{1cm}|p{1cm}|p{1cm}|p{1cm}|p{1cm}|p{1cm}|p{1cm}|p{1cm}|p{2.5cm}|}
				\hline
					Aufgabe & \centering 1 & \centering 2 & \centering 3 &  \centering 4 & \centering 5 & \centering 6 & \centering 7 & \centering 8 & Summe \\
				\hline
					Punkte &  &  &  &  &  &  &  &  & \\
				\hline
					Erreicht &  &  &  &  &  &  &  &  & \\
				\hline
			
		\end{tabularx}
		
		\vspace{0.5cm}
		
		Für eine Veröffentlichung der Note im Learnweb unterschreiben Sie bitte hier auf der Linie: \\
        
		\vspace{1cm}
        
		\begin{tabularx}{\linewidth}{Xr}\cline{2-2}
			&  Hildesheim, den 17.10.2025 
		\end{tabularx}

\newpage
%% -------
%% Front page end
%% -------

\thispagestyle{fancy} %% activate header/footer
\setcounter{page}{1} 

%% -------
%% Topic 1
%% -------
\section{This is topic 1}
Some description may go here. Can include text, images, source code, and other things. 

\subsection{5P}
This is the first task

\begin{solution}
    This is the solution for task a)
\end{solution}

\subsection{4P}
This is the second task

\begin{solution}
    This is the solution for task b)
\end{solution}


%% -------
%% Topic 2
%% -------
\section{This is topic 2}
Some description may go here. Can include text, images, source code, and other things. 

\subsection{2P}
This is the first task

\begin{solution}
    This is the solution for task a)
\end{solution}


\subsection{3P}
This is the second task

\begin{solution}
    This is the solution for task b)
\end{solution}


\subsection{4P}
This is the third task

\begin{solution}
    This is the solution for task c)
\end{solution}

\end{document}